% page 418 du fac-simile (image p-417 du scan)
% bloc 157.5 x 233.3 mm ; marge gauche 8.0 mm
\pgc{-12.700mm}{0.000mm}{
\l{\cel{3}410.}
\l{}
\l{\sou{oriento}. I. (\sou{literaturo}) Punto di la horizonto de ube suno}
\l{\cel{3}semblas levar su. - II. Regiono di la terglobo, este di}
\l{\cel{3}Europa (precipue Avan-Azia ed Egiptia). III. (\sou{framason}-}
\l{\cel{3}\sou{ismo}).\cel{1}\sou{Granda Oriento}\cel{1}: Quaza lojio centra, en la urbo}
\l{\cel{3}precipua di lando, qua konsistas ye la reprezenteri di}
\l{\cel{3}la lojii framasonala di la provinci.- DEFIRS.}
\l{}
\l{\sou{orifico}. Aperturo qua formacas la enireyo di kavajo. - EFIS.}
\l{}
\l{\sou{oriflamo}. Mikra standardo di qua la parto qua fluktuas finas}
\l{\cel{3}per du pinti. - DEFIS.}
\l{}
\l{\sou{origano}. (\sou{bot.}) Planto aromatoza, ek la familio \textquotedbl{}labiei\textquotedbl{},}
\l{\cel{3}majorano sovaja. - L.\cel{1}\sou{origanum}. - DEFIS.}
\l{}
\l{\sou{origino}. I. Departo-punto (loko, raso, e c.) di la nasko di}
\l{\cel{3}individuo, familio, populo - di la formaco di ulo. -}
\l{\cel{3}II. Instanto di la nasko. - III. (\sou{geom.}) Punto de qua}
\l{\cel{3}onu kontas la koordinati. - DEFIRS.}
\l{}
\l{\sou{orikalko}. (en\cel{1}\sou{la epoki antiqua}). Kompozuro metala, analoga}
\l{\cel{3}a latuno. - DEFIS.}
\l{}
\l{\sou{oriolo}.\cel{2}(\sou{zool.})\cel{2}Pasero di qua la plumaro esas flava. -}
\l{\cel{3}L.\cel{1}\sou{oriolus}.\cel{2}- EF.}
\l{}
\l{\sou{oripelo}. I. Filo o folio de latuno, polisita, qua, de fore,}
\l{\cel{3}brilas quale oro, e de qua onu facas ornamenti. - II.}
\l{\cel{3}Folieto dina, oratra, quan onu uzas por oro-sembligar}
\l{\cel{3}objekti diversa. - FI.}
\l{}
\l{\sou{orjato}. Siropo facita, olim, ek dekokturo de hordeo - e,}
\l{\cel{3}nun, ek emulsuro de mandeli. - DEFIS.}
\l{}
\l{\sou{orkestro}. I. (1) (en\cel{1}\sou{la epoki antiqua)}\cel{1}. Parto di la teatro}
\l{\cel{3}Greka avan e sub la ceneyo (\textquotedbl{}stejo) ube la koro stacis}
\l{\cel{3}ed evolucionis. (2) Parto korespondanta di la teatro mo-}
\l{\cel{3}derna, ube esas la muzikisti qui pleas muzikaji koncertala.}
\l{\cel{3}- II. Ensemblo de ta muzikisti.- DEFIRS.}
\l{}
\l{\sou{orkideo}. (\sou{bot.})\cel{2}Un ek la familii de planti monokotiledona,}
\l{\cel{3}bulboza, de qua la \textquotedbl{}orkis\textquotedbl{} esas la tipo. L.\cel{1}\sou{orchidea}.}
\l{\cel{3}- DEFIRS.}
\l{}
\l{\sou{orlar}. (\sou{trans.})\cel{2}Facar la bordo di texuro, per rifaldo e}
\l{\cel{3}suto, por impedar la destexeso.- FIS.}
\l{}
\l{\sou{ornar}.\cel{2}(\sou{trans.})\cel{2}Kovrar per vesti eleganta, per ulo qua}
\l{\cel{3}plubeligas, per ornamenti. -\cel{1}\sou{Anke metaf}. - EFIS.}
\l{}
\l{ornamento. Detalo adaptita ad ensemblo por igar lu plu}
\l{\cel{3}bela.- DEFIS.}
}
